
\section{Introdução}

\par Para a disciplina de \textbf{Projeto de Software} seu objetivo é analisar as metodologias de desenvolvimento/gerenciamento de projetos ágil, visando uma melhora no uso do tempo de desenvolvimento e a própria agilidade no desenvolvimento do projeto.
\par De forma sucinta o movimento ágil é uma alternativa a gestão tradicional de projetos, as abordagem ágis nos tempos atuais vem sendo se suma importância e um divisor de águas para que um projeto ou \textit{software} se manter no mercado, pois cada vez esta um ambiente incerto e turbuleno.
\par Segundo \citeonline{smartsheet}, os quatro valores do manifesto ágil são:

 \begin{enumerate}[label=\Roman{*}, ref=(\roman{*})]
   \item Indivíduos e interações mais que processos e ferramentas
   \item Software em funcionamento mais que documentação abrangente
   \item Colaboração com o cliente mais que negociação de contratos
   \item Responder a mudanças mais que seguir um plano
 \end{enumerate}
 \par E desta maneira possuem os \textbf{doze} princípios do manifesto ágil, e que são:
\begin{asparaenum}
  \item Satisfação do cliente por meio do fornecimento contínuo e adiantado de software.
  \item Acomodação de mudanças de requisito durante todo o processo de desenvolvimento.
  \item Entrega frequente de software em funcionamento.
  \item Colaboração entre as partes interessadas do negócio e os desenvolvedores em todo o projeto.
  \item Apoio, confiança e motivação das pessoas envolvidas.
  \item Possibilidade de interações presenciais.
  \item Software em funcionamento é a principal medida de progresso.
  \item Processos ágeis para dar apoio a um ritmo de desenvolvimento consistente.
  \item Atenção aos detalhes técnicos e ao design aumenta a agilidade.
  \item Simplicidade.
  \item Equipes auto-organizáveis realizam ótimas arquiteturas, requisitos e designs.
  \item Reflexões periódicas sobre como aumentar a eficácia.

  \cite{smartsheet}
\end{asparaenum}

\par Esta aula prática tem por objetivo o desenvolvimento das etapa de um projeto hipotetico ágil, com o fim de aplicar os conhcimentos adqueridos na cadeira de \textbf{projeto de \textit{software}}.

\section{Métodos}
\par Para o desenvovimento desta aula prática foi disponibilizado um roteiro, e segindo este roteiro é porposto elaboração em duas etapas, sendo elas:
\begin{enumerate}
  \item Se colocar no lugar do cliente, e pensar em um aplicativo para a confecção.
  \item E se colocar no lugar do \textit{Product Owner}, que deverá elaborar o aplicativo proposto.
\end{enumerate}


\noindent \begin{minipage}[c]{0.6\textwidth}
  \vspace {1cm}
\par Após as analise iniciais, o roteiro propoe o uso de algumas ferramentas para controle de etapas \textbf{(IceScrumm, Trello e Asna)}, opto pelo uso do \href{https://trello.com/home}{Trello}, esta ferramenta proporciona a integração com o \href{www.github.com}{GitHub}, adicionado e comentando ramificações e solicitações de mesclagem de fragmentos de códigos.

\end{minipage}
\begin{minipage}[c]{0.4\textwidth}
  \captionof{figure}{Logo Trello}
  \includegraphics[width=\textwidth]{figure/trello.png}
  \label{fig:log_trello}
  {\fontsize{10pt}{\baselineskip} \selectfont Fonte: \citeonline{logtrello}}
\end{minipage}




\section{Resultados}
\par A eleboração desta aula prática esta dividida em duas partes, a primeira, no qual serei o cliente e deverá propor um aplicativo que se deseja desenvolver, levantar as funcionalidades e caracteristicas. Na segunda etapa, deverá ser o \textit{Product Owner}, quem ira desenvolver o aplicativo proposto e as responsabilidades desta etapa são:
\begin{itemize}
  \item Definir as funcionalidades do produto, ou seja, desenvolver o product backlog;
  \item Priorizar as funcionalidades de acordo com o valor de negócio;
  \item Montar um quadro do Scrum (Kanban) com as divisões de etapas, tarefas, data de entrega e responsáveis por atividade. Para este item, imagine que o desenvolvimento do seu aplicativo está em um estágio mais avançado, por este motivo, deve haver tarefas em todas as etapas. Utilize uma das ferramentas propostas para montar o seu quadro.
\end{itemize}

\subsection{Primeira etapa}


\subsection{Segunda etapa}




$\underline{{\color{Blue} }}{\color{red} \, \: \: \left ( \iint_{\widetilde{ \theta }}^{\phi } \right )}  $



\par Aqui vai um exemplo de código em \LaTeX2e.


\begin{figure}[H] %Figuras da aula pratica 1.1
  \center
  \caption{Resultado da aula prática 1.2}
  \subfigure[ Algoritmo.\label{fig:pri2}]{\includegraphics[scale=0.4]{figure/placeholder.jpg}}
  \subfigure[Comportamento.\label{fig:seg2}]{\includegraphics[scale=.4]{figure/placeholder.jpg}}
  % Deixar o espço

  {\fontsize{10pt}{\baselineskip}\selectfont Fonte: \citeonline{alguem2022nada}}
  \label{fig:ap1_cod_vigual1}
\end{figure}



Para referenciar utilize \cite{ninguem2022curioso}. Também pode ser citado integrada ao texto, de acordo com \citeonline{alguem2022nada}.

Para inserir imagens adicione a figura no diretório \textit{/figure}


\par Estou usando \href {https://cocalc.com/} {CoCal}

E para referenciar a figura \ref{fig:imagem_massa} utilize dessa forma.

\begin{equation}
  a + b = c \label{1}
  $$ $$ %usado para pular linha
  \sqrt{x} + \sqrt{\smash[b]{y}} + \sqrt{z}
\end{equation}


\par Exemplo de incerção de formula \eqref{1},


Aqui é um exemplo de rodapé. \footnote{Um exemplo de rodapé}


\begin{enumerate}[label=\Roman{*}, ref=(\roman{*})]
  \item fsfsdf
  \item kugfhiuh
\end{enumerate}

\begin{asparaenum}
\item Anterior ... \cite{ninguem2022curioso}
\item Próximo ... \label{pl1}
\end{asparaenum}


\section{Conclusões}

\par Podemos aferir com elaboração da aula prática, a necessidade do uso de metodologias com na gestão de software.



  %$X \xLongleftarrow[\text{NATAN}]{\text{OGLIARI}} Y $ %COM TEXTO
	% $\uparrow$ %Seta para Cima
	%$\overleftarrow{NATAN}$
